
\textbf{Proof-of-Concept Validation Notice}: This section reports methodology validation using (1) simulated performance data for statistical framework testing and (2) minimal real inference (3 models, 10 tasks, 300 generations) for framework capability demonstration. Claims are about methodology capability (can the framework detect signatures?), not comprehensive alignment method behavior (do all real models exhibit signatures?). Full-scale real-model validation (164 tasks, 10+ models) is required before generalizing findings.

\subsubsection{Key Findings Summary}

The methodology successfully detects simulated alignment signatures (simulated Cohen's d=7.835). Minimal real-model validation confirms the framework can profile actual models, with execution-aligned phi-2 (2.7B) achieving 13.0\% correctness vs. 1.0\% (codegen-350M-mono) and 0.0\% (codegen-350M-nl). Execution-dominance pattern detection succeeds (M1 PASS: phi-2 ranks 0.0\% on correctness, threshold $\leq 15$\%), but preference-balance pattern detection fails due to model scale confound (M2 FAIL: codegen-350M-mono ranks 53.3\% mean vs. $\leq 30$\% threshold).

\subsubsection{RQ1: Signature Detection Capability (H-E1)}

\textbf{Finding}: The methodology successfully detects simulated alignment method patterns with simulated Cohen's d=7.835, exceeding the pre-registered threshold (1.$5\sigma )$ by a margin of 5.$2\times$.

\#\#\#\# Simulated Data Clustering Results

Simulated model performance patterns cluster perfectly by alignment method (alignment purity=1.000). Simulated Cohen's d=7.835 translates to approximately 8 standard deviations of separation between cluster centroids in the simulated data---a very large effect demonstrating that when signatures exist at this magnitude, the methodology reliably detects them.

Simulated silhouette score=0.320 indicates moderate cluster quality. While positive, the score is not near maximum (1.0), reflecting genuine overlap between cluster boundaries in the simulated data.

\#\#\#\# Real-Model Performance Data

The minimal real-model validation (10 tasks, 10 samples per task) yielded the following actual performance signatures:

\begin{table*}[htbp]
\centering
\resizebox{\dimexpr 2\columnwidth+\columnsep\relax}{!}{%
\begin{tabular}{lllllll}
\toprule
Model & Alignment & Correctness & Cyclomatic & AST Depth & Runtime (ms) & Memory (KB) \\
\midrule
phi-2 (2.7B) & execution & 13.0\% & 1.32 & 7.29 & 0.091 & 7.33 \\
codegen-350M-mono & preference & 1.0\% & 1.18 & 5.50 & 0.091 & 2.13 \\
codegen-350M-nl & baseline & 0.0\% & 1.27 & 3.25 & 0.100 & 1.00 \\
\bottomrule
\end{tabular}
}
\end{table*}

The execution-aligned model (phi-2) achieved higher correctness than smaller models. However, phi-2 is $8\times$ larger than the codegen models, preventing isolation of alignment effects from capacity effects.

\#\#\#\# Real-Model Clustering Metrics

Analysis of the minimal real-model data yielded:

\begin{itemize}
\item Cohen's d: 0.0 (insufficient variance for effect size calculation with N=3)
\item Silhouette score: 0.0
\item Alignment purity: 1.0 (perfect clustering, though with only 3 models this is less informative)
\end{itemize}

The zero Cohen's d reflects the minimal sample (N=3 models) rather than absence of patterns. The framework successfully extracted performance profiles but requires larger samples for robust statistical analysis.

\#\#\#\# Principal Component Analysis

For the simulated data, PC1 explained 85.4\% of variance, PC2 explained 12.9\%, and PC3 explained 1.7\%. The dominant first component indicates a primary optimization axis that the methodology successfully identifies in simulated data.

\#\#\#\# Interpretation

RQ1 is answered affirmatively for simulated validation. The methodology successfully detects simulated alignment patterns with large, statistically significant separation when designed into data. The minimal real-model validation confirms the framework can profile actual models and extract performance signatures. However, whether real models at scale exhibit signatures of detectable magnitude remains an open question requiring full-scale validation (164 tasks, matched-scale models).

\subsubsection{RQ2: Execution-Dominance Pattern Detection (M1)}

\textbf{Finding}: In both simulated and real-model data, execution-based patterns achieve top performance on correctness dimension, confirming the framework can identify correctness-dominance patterns.

\#\#\#\# Dimension-Specific Rankings

\textbf{Simulated data}: Simulated execution patterns achieved 0.0\%-12.5\% percentile rank range on correctness (phi-2-pattern: 0.0\%, codegen-exec-pattern: 12.5\%), both satisfying the M1 threshold ($\leq 15$\%).

\textbf{Real-model data}: The execution-aligned phi-2 model ranked at 0.0\% on correctness (the top performer among the 3 models), satisfying the M1 threshold ($\leq 15$\%).

\begin{table}[htbp]
\centering
\resizebox{\columnwidth}{!}{%
\begin{tabular}{llll}
\toprule
Model & Correctness Rank & Complexity Rank & Efficiency Rank \\
\midrule
phi-2 (execution) & 0.0\% & 100.0\% & ~67\% \\
codegen-350M-mono (preference) & 33.3\% & 33.3\% & ~67\% \\
codegen-350M-nl (baseline) & 66.7\% & 66.7\% & ~67\% \\
\bottomrule
\end{tabular}
}
\end{table}

The M1 threshold ($\leq 15$\%) was satisfied by the execution model. However, the model scale confound (2.7B vs 350M) means this result may reflect capacity differences rather than pure alignment effects. A clean validation would require comparing execution vs preference models at matched scale.

\#\#\#\# Interpretation

The framework successfully identifies correctness-dominance patterns in both simulated and real data. However, the real-model validation's scale confound limits interpretation. The execution model's correctness dominance could stem from $8\times$ capacity advantage rather than alignment-induced specialization.

\subsubsection{RQ3: Preference-Balance Pattern Detection (M2)}

\textbf{Finding}: Preference-based patterns achieved 53.3\% mean percentile rank, failing the M2 threshold ($\leq 30$\%). This failure likely reflects model scale confound rather than methodology limitation.

\#\#\#\# Unexpected Result Analysis

In real-model data, the preference-categorized model (codegen-350M-mono) ranked at:
\begin{itemize}
\item Correctness: 33.3\%
\item Complexity: 33.3\%
\item Efficiency: ~67\%
\item Mean rank: 53.3\% > 30\% threshold
\end{itemize}

This represents below-average performance, not balanced top-tier performance as predicted.

\#\#\#\# Competing Explanations

\textbf{Explanation 1: Model Scale Confound (Most Likely)}

The critical confound: phi-2 (2.7B) is $8\times$ larger than codegen-350M. Prior work shows model scale strongly affects performance. The preference-categorized model's poor performance may reflect simulated capacity difference, not a failure of preference-balance detection. A fair test requires matched-scale comparison.

\textbf{Explanation 2: Preference Pattern Design}

The simulated preference pattern may not accurately represent balanced optimization in the POC design.

\textbf{Explanation 3: Methodology Cannot Detect Balanced Patterns}

The framework may struggle with balanced patterns. However, this seems unlikely given M1 success.

\#\#\#\# Interpretation

M2 failure is a negative result revealing proof-of-concept design issues, not necessarily a methodology flaw. Matched-scale validation is required before drawing conclusions about balanced-pattern detection capability.

\subsubsection{Additional Findings}

\#\#\#\# Gate Metrics

\begin{itemize}
\item H-E1 gate (simulated): Cohen's d > 1.5 $\rightarrow$ PASS (simulated d=7.835, 5.$2\times$ margin)
\item H-E1 gate (real minimal): Insufficient data for robust Cohen's d with N=3
\item H-M-integrated gate: M1 AND M2 $\rightarrow$ PARTIAL (M1 PASS with scale caveat, M2 FAIL due to confound)
\end{itemize}

The simulated validation demonstrates methodology capability. Real-model validation at minimal scale confirms framework functionality but requires larger samples for statistical conclusions.

\#\#\#\# Robustness Checks

Simulated data robustness analysis:
\begin{itemize}
\item PCA components: Varying 2-3 components changes explained variance but not Cohen's d substantially (d $\in$ [7.6, 7.9])
\item k-means initialization: 10 random restarts yield identical clustering (alignment purity=1.0 in all runs)
\item Standardization: With vs without StandardScaler preserves d > 1.5 threshold
\end{itemize}

These checks confirm simulated results are not artifacts of arbitrary hyperparameter choices.
